%! TEX root = main.tex  % Specifies the main document file for multi-file projects

\chapter{Course Introduction}

The course Computer Music: Languages and Systems (Course number: 054283) is designed to provide students with a comprehensive understanding of music programming and digital audio systems. It offers both theoretical lectures and practical laboratory sessions, ensuring a hands-on approach to learning.

\section{Course Organization}

The course is divided into three main parts:

\begin{enumerate}
    \item SuperCollider as a music programming language (8 hours lecture, 6 hours lab, 2 hours Q/A and exam preparation).
    
    \begin{itemize}
        \item Architecture of SuperCollider
        \item Coding in SuperCollider
    \end{itemize}
    
    \item Computer Music Systems and plugin architecture (6 hours lecture, 6 hours lab, 2 hours Q/A and exam preparation).
    
    \begin{itemize}
        \item Digital Audio Workstations (DAWs)
        \item Architecture and development of plugins using the JUCE platform
    \end{itemize}

    \item Interaction Design (8 hours lecture, 8 hours lab, 4 hours Q/A and exam preparation).
    
    \begin{itemize}
        \item Machine-to-machine (M2M) and human-to-machine (H2M) interaction
        \item M2M: MIDI and OSC protocols
        \item H2M: Arduino and Bela platforms
    \end{itemize}
\end{enumerate}

A detailed schedule is available on WeBeep.

\section{Textbooks and Materials}

Since computer music is a broad research field, no single textbook covers all topics of the course. Instead, students will rely on:

\begin{itemize}
    \item Course slides, which provide comprehensive coverage.
    \item References at the end of each slide set, some of which are available online.
\end{itemize}

\section{Laboratories}

Laboratory sessions will be held in lecture rooms. Students must install the necessary software on their own laptops, including:

\begin{itemize}
    \item \href{http://supercollider.github.io/download}{SuperCollider}
    \item \href{https://juce.com}{JUCE}
    \item An IDE of choice (Xcode, Visual Studio, etc.)
    \item Arduino, Processing, and other required software
\end{itemize}

\section{Exam and Homework}

The assessment consists of:

\begin{itemize}
    \item A homework assignment (max 14 points), where students develop:
    \begin{itemize}
        \item A synth/effect in SuperCollider
        \item A plugin in JUCE
        \item Interaction via MIDI/OSC and hardware interfaces (Arduino, Bela, etc.)
    \end{itemize}
    \item A written exam (max 18 points), which includes:
    \begin{itemize}
        \item Theoretical questions
        \item Practical exercises on SuperCollider, JUCE, and interaction design
    \end{itemize}
\end{itemize}

Homework is done in groups (3–5 students). A GitHub repository will be used for submission. The final presentation will be graded based on implementation quality, presentation skills, and creativity.